%==============================================================
% APUNTES DE DERECHO CONSTITUCIONAL PROFUNDIZADO
% Documento autónomo (sin plantillas ni archivos externos)
%==============================================================
\documentclass[12pt,oneside]{book}

%--------------------------------------------------------------
% METADATOS
%--------------------------------------------------------------
\newcommand{\universidad}{Universidad de Buenos Aires (UBA)}
\newcommand{\facultad}{Facultad de Derecho}
\newcommand{\materia}{Derecho Constitucional Profundizado}
\newcommand{\titulo}{Apuntes de Derecho Constitucional Profundizado}
\newcommand{\subtitulo}{Interpretación constitucional, control de constitucionalidad y precedentes}
\newcommand{\autor}{Isaac Edgar Camacho Ocampo}
\newcommand{\anio}{2025}

%--------------------------------------------------------------
% PAQUETES
%--------------------------------------------------------------
\usepackage[utf8]{inputenc}
\usepackage[T1]{fontenc}
\usepackage[provide=*,spanish]{babel}

\usepackage[
    a4paper,
    top=2.5cm,
    bottom=2.5cm,
    left=3cm,
    right=2.5cm,
    headheight=15pt
]{geometry}

\usepackage{amsmath}
\usepackage{amssymb}
\usepackage{mathtools}

\usepackage{graphicx}
\usepackage{float}
\usepackage{caption}
\usepackage{subcaption}

\usepackage{booktabs}
\usepackage{array}
\usepackage{longtable}
\usepackage{tabularx}

\usepackage{enumitem}

\usepackage{xcolor}
\usepackage[most]{tcolorbox}

\usepackage{fancyhdr}
\usepackage{ragged2e}

\usepackage[
    colorlinks=true,
    linkcolor=blue!50!black,
    citecolor=green!40!black,
    urlcolor=blue!60!black,
    pdftitle={\titulo},
    pdfauthor={\autor},
    pdfsubject={Apuntes universitarios},
    bookmarks=true
]{hyperref}

%--------------------------------------------------------------
% CONFIGURACIÓN
%--------------------------------------------------------------
\graphicspath{
    {./img/}
    {./graficos/}
    {./figuras/}
}

\setcounter{tocdepth}{2}

\setlength{\parindent}{0pt}
\setlength{\parskip}{0.5em}

\setlist[itemize]{
    topsep=0.3em,
    itemsep=0.2em
}

\setlist[enumerate]{
    topsep=0.3em,
    itemsep=0.2em
}

\clubpenalty=10000
\widowpenalty=10000

%--------------------------------------------------------------
% COMANDOS
%--------------------------------------------------------------
\newcommand{\abs}[1]{\left|#1\right|}
\newcommand{\norm}[1]{\left\lVert#1\right\rVert}
\newcommand{\vect}[1]{\mathbf{#1}}
\newcommand{\deriv}[2]{\frac{d#1}{d#2}}
\newcommand{\pderiv}[2]{\frac{\partial#1}{\partial#2}}

%--------------------------------------------------------------
% ENTORNOS / CAJAS
%--------------------------------------------------------------
\newtcolorbox{definition}[1]{
    enhanced,
    breakable,
    title={Definición: #1},
    fonttitle=\bfseries,
    colback=blue!3,
    colframe=blue!50!black,
    boxrule=0.8pt,
    arc=2mm
}

\newtcolorbox{important}{
    enhanced,
    breakable,
    title={Importante},
    fonttitle=\bfseries,
    colback=yellow!8,
    colframe=orange!70!black,
    boxrule=0.8pt,
    arc=2mm
}

\newtcolorbox{example}[1]{
    enhanced,
    breakable,
    title={Ejemplo: #1},
    fonttitle=\bfseries,
    colback=green!4,
    colframe=green!40!black,
    boxrule=0.8pt,
    arc=2mm
}

\newtcolorbox{formula}[1]{
    enhanced,
    breakable,
    title={Fórmula / Regla: #1},
    fonttitle=\bfseries,
    colback=purple!4,
    colframe=purple!50!black,
    boxrule=0.8pt,
    arc=2mm
}

\newtcolorbox{exercise}[1]{
    enhanced,
    breakable,
    title={Ejercicio: #1},
    fonttitle=\bfseries,
    colback=gray!5,
    colframe=gray!60!black,
    boxrule=0.8pt,
    arc=2mm
}

\newtcolorbox{note}{
    enhanced,
    breakable,
    title={Nota},
    fonttitle=\bfseries,
    colback=white,
    colframe=gray!50,
    boxrule=0.6pt,
    arc=2mm
}

%--------------------------------------------------------------
% CONFIGURACIÓN FINAL
%--------------------------------------------------------------
\pagestyle{fancy}
\fancyhf{}
\fancyhead[L]{\small\nouppercase{\leftmark}}
\fancyhead[R]{\small\materia}
\fancyfoot[C]{\thepage}
\renewcommand{\headrulewidth}{0.4pt}
\renewcommand{\footrulewidth}{0pt}

\fancypagestyle{plain}{
    \fancyhf{}
    \fancyfoot[C]{\thepage}
    \renewcommand{\headrulewidth}{0pt}
}

%==============================================================
\begin{document}
%==============================================================

%--------------------------------------------------------------
% PORTADA
%--------------------------------------------------------------
\begin{titlepage}

\thispagestyle{empty}

\begin{center}

\vspace*{1cm}

{\large \textsc{\universidad}}

\vspace{0.2cm}

{\large \textsc{\facultad}}

\vspace{3cm}

{\Huge \textbf{\materia}}

\vspace{1cm}

{\LARGE \subtitulo}

\vspace{0.8cm}

\textit{Estudiar con constancia es la mejor manera de convertir los conceptos en criterio propio.}

\vspace{1.2cm}

\rule{0.70\textwidth}{0.5pt}

\vspace{0.5cm}

{\large Apuntes de estudio}

\vspace{0.5cm}

\rule{0.70\textwidth}{0.5pt}

\vfill

{\large \autor}

\vspace{0.3cm}

{\large \anio}

\end{center}

\vfill

\begin{flushleft}
\textit{Este es un modesto aporte a los estudiantes de ``Derecho Constitucional Profundizado'' no pretende sustituir el material oficial de la materia.}
\end{flushleft}

\end{titlepage}

%--------------------------------------------------------------
% ÍNDICE
%--------------------------------------------------------------
\tableofcontents

%--------------------------------------------------------------
% PRESENTACIÓN
%--------------------------------------------------------------
\chapter*{Presentación de los contenidos}
\addcontentsline{toc}{chapter}{Presentación de los contenidos}
\markboth{Presentación de los contenidos}{Presentación de los contenidos}

Los contenidos se organizan en tres bloques que forman una misma cadena de razonamiento: \textbf{interpretación} $\rightarrow$ \textbf{control} $\rightarrow$ \textbf{precedente}.

Ante un caso en el que se cuestiona la constitucionalidad de una ley, es necesario, en primer lugar, saber cómo leer e interpretar las normas (interpretación jurídica y constitucional). En segundo lugar, es necesario conocer cuál es el poder del juez para declarar inconstitucional esa ley y bajo qué sistema de control opera (control de constitucionalidad). Finalmente, corresponde determinar si existe un fallo anterior de la Corte que ya resolvió la cuestión y si el juez está obligado a seguirlo (precedentes constitucionales).

\begin{note}
\textbf{Relevancia para el ejercicio profesional.} Los temas abordados constituyen el núcleo del ejercicio profesional del abogado y de cualquier operador jurídico. Comprender la interpretación constitucional permite argumentar con solidez ante los tribunales, redactar escritos con fundamento y anticipar cómo los jueces resolverán los conflictos. El control de constitucionalidad difuso, propio del sistema argentino, significa que cualquier abogado, en cualquier fuero y en cualquier instancia, puede plantear la invalidez de una norma. Los precedentes de la Corte Suprema organizan la práctica diaria: saber qué fallos obligan a los jueces inferiores evita litigios innecesarios y permite asesorar a los clientes con mayor certeza.
\end{note}

%==============================================================
% CAPÍTULO 1
%==============================================================
\chapter{Interpretación jurídica y constitucional}
\markboth{Capítulo \thechapter. Interpretación}{}

\section{La letra como primera fuente de interpretación}

La interpretación jurídica y la interpretación constitucional propiamente dicha se apoyan en un conjunto de técnicas de hermenéutica que se basan, esencialmente, en el análisis concreto de la letra de la Constitución y de las normas inferiores.

La \textbf{letra de la Constitución} y la letra de las normas que deban interpretarse constituyen la primera y esencial fuente de interpretación. Como toda hermenéutica, esta tarea puede requerir además una interpretación conjunta de dos normas: por ejemplo, cuando hay que interpretar dos normas que puedan estar en pugna o que ofrezcan algún tipo de contradicción, para determinar cuál prevalece en el caso y cuál resulta aplicable.

Resulta ineludible considerar que la letra de la Constitución es la primera, básica y primordial fuente de interpretación.

\section{El textualismo}

\subsection{Concepto y características}

\begin{definition}{Textualismo}
Corriente que afianza el valor de la \textbf{letra como primera fuente} de interpretación y le da un sentido estricto: la letra no solo es la primera referencia a la hora de interpretar el cuerpo legal que se quiere aplicar, sino que además la interpretación, o la creatividad, del juez u órgano que realiza la aplicación queda muy reducida, porque es el texto el que marca los límites.
\end{definition}

De este modo, el textualismo reduce el margen creativo del juez: las palabras de la norma marcan los límites de la interpretación posible.

\subsection{El textualismo en la jurisprudencia argentina}

En la jurisprudencia argentina, la Corte sostiene, respecto de esta fuente textualista, que la primera fuente de interpretación es la letra de la norma y que esta debe entenderse según el \textbf{uso corriente o significado corriente} que las palabras tienen en los tiempos actuales.

Se trata de una modulación más flexible que la del textualismo puro: la Corte argentina no congela el significado de las palabras en el momento histórico de la sanción, sino que las interpreta con el sentido que tienen en la actualidad.

\subsection{Aplicación práctica: artículo 19 de la Constitución Nacional}

El artículo 19 de la Constitución Nacional consagra el \textbf{principio de reserva}:

\begin{quote}
\textit{Las acciones privadas de los hombres que de ningún modo ofendan al orden y a la moral pública, ni perjudiquen a un tercero, están sólo reservadas a Dios y exentas de la autoridad de los magistrados. Ningún habitante de la Nación será obligado a hacer lo que no manda la ley, ni privado de lo que ella no prohíbe.}
\end{quote}

Al analizar el principio de reserva y determinar si las acciones privadas ofenden o no las buenas costumbres, y si una determinada actividad se encuentra o no en la zona de reserva, corresponde apelar al significado que las palabras tienen actualmente y de forma corriente para desentrañar la interpretación constitucional.

\begin{example}{Tenencia de estupefacientes para consumo personal y artículo 19}
La tenencia de estupefacientes para consumo personal plantea la cuestión de cómo contrastarla con el principio de reserva del artículo 19 de la Constitución Nacional. Surgen entonces interrogantes: qué extensión darle a las expresiones según las cuales las acciones no pueden ofender la moral ni las buenas costumbres, ni perjudicar a terceros; cómo interpretar las \textit{buenas costumbres} y cómo desentrañar su alcance; y cómo debe interpretarse en los tiempos actuales todo el bagaje que tiene la redacción de la norma.

El ejemplo refleja las distintas posiciones posibles: un ámbito de interpretación más acotado, en el que el intérprete se guía por la literalidad del texto, o bien un margen para analizar más, adaptar y proyectar la norma a los tiempos actuales, entendiendo que el cambio de circunstancias puede generar un cambio en la postura y en el alcance que se le da a las normas.
\end{example}

\section{El originalismo}

\subsection{Concepto y relación con el textualismo}

El originalismo \textbf{no es equivalente} al textualismo. Ambas son corrientes que intentan ser restrictivas de la actividad de los jueces, pero por vías distintas:

\begin{itemize}
    \item \textbf{Textualismo:} la letra es la primera fuente de interpretación y el juez debe limitarse a leer con fidelidad los términos y a atender a su uso normal y corriente para descifrar el sentido y la aplicación de la norma.
    \item \textbf{Originalismo:} a la idea de que la letra es la fuente primaria le agrega un sentido y un peso histórico.
\end{itemize}

\begin{definition}{Originalismo}
Corriente según la cual toda interpretación legal, si bien debe responder a la fidelidad del texto que describe la situación, debe además estar imbuida y consustanciada con el \textbf{período histórico} en el que la norma fue redactada. Se afianza la redacción como fuente original de la interpretación, pero consustanciada con el momento histórico en que fue puesta en vigencia; por lo tanto, las palabras tienen el sentido que tenían históricamente al momento de su redacción y sanción.
\end{definition}

\subsection{Comparación entre textualismo y originalismo}

\begin{table}[H]
\centering
\small
\begin{tabularx}{\textwidth}{
    >{\raggedright\arraybackslash}p{0.24\textwidth}
    >{\raggedright\arraybackslash}X
    >{\raggedright\arraybackslash}X
}
\toprule
\textbf{Dimensión} & \textbf{Textualismo} & \textbf{Originalismo} \\
\midrule
\textbf{Fuente primaria} & Letra de la norma & Letra y contexto histórico de sanción \\
\textbf{Significado de las palabras} & Uso corriente y actual & Sentido al momento de la redacción \\
\textbf{Rol del juez} & Muy limitado por el texto & Limitado por texto e historia legislativa \\
\textbf{Referencia cultural} & El texto en sí mismo & Voluntad de los padres fundadores \\
\textbf{Adaptación temporal} & Palabras con sentido actual & Sin adaptación: anclaje histórico \\
\textbf{Arraigo en Argentina} & Moderado (doctrina de la Corte) & Débil (más fuerte en EE.UU.) \\
\textbf{Arraigo en EE.UU.} & Fuerte & Muy fuerte (jueces conservadores) \\
\bottomrule
\end{tabularx}
\end{table}

\subsection{El originalismo en los Estados Unidos}

En el derecho argentino el originalismo no está tan arraigado como en los Estados Unidos. En ese país las corrientes originalistas y tradicionalistas están más afianzadas, porque la labor de los jueces genera la idea de que hay que mantenerse muy ceñido al contexto histórico y a la voluntad de los legisladores.

Un punto clave es el concepto de los \textbf{padres fundadores}, que en la cultura americana tiene un desarrollo que no existe con la misma intensidad en la argentina. Las fuentes jurídicas americanas expresan una deferencia muy explícita hacia la idea de interpretar lo más fielmente posible la voluntad de esos legisladores que fueron los redactores de la Constitución y de las leyes esenciales.

Esta situación tiene un condicionante práctico importante: en la Argentina no existe constancia documentada clara de cómo fue la Asamblea Constituyente del año 1853. La de 1860 está más documentada, lo que explica en parte por qué no existe la misma tradición de apelación a los constituyentes originarios que sí existe en los Estados Unidos.

Al analizar la Corte americana puede advertirse que los jueces más tradicionales y conservadores son quienes adhieren a la corriente originalista. Esto se plasma en la interpretación jurídica: los operadores del derecho originalistas tratan de pensar cómo responderían ante la situación actual figuras como Franklin o Madison y cómo aplicarían la Constitución en la actualidad.

\begin{important}
Tanto el textualismo como el originalismo son corrientes conservadoras desde el punto de vista metodológico: priorizan la fidelidad de la letra y del momento histórico por sobre cualquier adaptación o proyección hacia los tiempos actuales. El efecto práctico es achicar el campo de interpretación de los jueces.
\end{important}

\section{Cuadro Cornell: interpretación jurídica y constitucional}

\begin{longtable}{@{}>{\RaggedRight\arraybackslash}p{0.27\textwidth}>{\RaggedRight\arraybackslash}p{0.68\textwidth}@{}}
\toprule
\textbf{Preguntas / palabras clave} & \textbf{Notas / respuestas} \\
\midrule
\endfirsthead
\toprule
\textbf{Preguntas / palabras clave} & \textbf{Notas / respuestas} \\
\midrule
\endhead
\bottomrule
\endlastfoot

\textbf{¿Cuál es la primera fuente de interpretación constitucional?} &
La letra de la Constitución y la letra de las normas que deban interpretarse, que constituyen la primera, básica y primordial fuente de interpretación. Puede requerir una interpretación conjunta de dos normas en pugna para determinar cuál prevalece en el caso. \\
\addlinespace[0.3em]\cmidrule[0.2pt]{1-2}\addlinespace[0.3em]

\textbf{¿Qué es el textualismo y cómo limita al juez?} &
Es la corriente que afianza el valor de la letra como primera fuente y le da un sentido estricto: la creatividad del juez queda muy reducida porque el texto marca los límites de la interpretación posible. \\
\addlinespace[0.3em]\cmidrule[0.2pt]{1-2}\addlinespace[0.3em]

\textbf{¿Cómo entiende la Corte argentina el significado de las palabras de la norma?} &
Con el uso o significado corriente que las palabras tienen en los tiempos actuales. Es una modulación más flexible que el textualismo puro, ya que no se congela el significado en el momento de la sanción. \\
\addlinespace[0.3em]\cmidrule[0.2pt]{1-2}\addlinespace[0.3em]

\textbf{¿En qué se diferencia el originalismo del textualismo?} &
Ambas son corrientes restrictivas de la actividad de los jueces. El textualismo se limita a la letra y a su uso corriente; el originalismo agrega el peso histórico: las palabras tienen el sentido que tenían al momento de su redacción y sanción, en consonancia con la voluntad de los padres fundadores. \\
\addlinespace[0.3em]\cmidrule[0.2pt]{1-2}\addlinespace[0.3em]

\textbf{¿Por qué está más arraigado el originalismo en los Estados Unidos que en la Argentina?} &
Porque en los Estados Unidos existe una deferencia explícita hacia la voluntad de los padres fundadores y los jueces conservadores adhieren a esta corriente. En la Argentina no existe constancia documentada clara de cómo fue la Asamblea Constituyente de 1853 (la de 1860 está más documentada). \\
\addlinespace[0.3em]\cmidrule[0.2pt]{1-2}\addlinespace[0.3em]

\textbf{¿Por qué el artículo 19 de la Constitución Nacional sirve como ejemplo de interpretación?} &
Porque obliga a decidir cómo interpretar expresiones como las buenas costumbres, la moral pública y el perjuicio a terceros. En el caso de la tenencia de estupefacientes para consumo personal, la respuesta depende de si se adopta una interpretación acotada a la literalidad o una que adapte la norma a los tiempos actuales. \\
\addlinespace[0.3em]\cmidrule[0.2pt]{1-2}\addlinespace[0.3em]

\midrule
\multicolumn{2}{@{}p{0.95\textwidth}@{}}{%
\textbf{Resumen:} La letra de la norma es la primera fuente de interpretación. El textualismo restringe al juez al texto, entendido en la jurisprudencia argentina con su significado corriente actual; el originalismo agrega el anclaje histórico al momento de la sanción y la voluntad de los padres fundadores, y tiene mayor arraigo en los Estados Unidos. Ambas corrientes son metodológicamente conservadoras y achican el campo de interpretación de los jueces.} \\
\end{longtable}

%==============================================================
% CAPÍTULO 2
%==============================================================
\chapter{Control de constitucionalidad}
\markboth{Capítulo \thechapter. Control de constitucionalidad}{}

\section{Conceptualización del control de constitucionalidad}

Toda interpretación constitucional supone desentrañar el sentido de la Constitución y determinar si existe una eventual colisión entre normas inferiores y la Constitución, o cuál es el propio alcance de la norma fundacional. Al hacerlo, se está ejerciendo un \textbf{control de constitucionalidad}.

\section{El sistema difuso argentino}

\subsection{Definición del control difuso}

El control de constitucionalidad argentino ha seguido muy fielmente la vertiente norteamericana.

\begin{definition}{Control de constitucionalidad difuso}
Sistema en el que el control de constitucionalidad \textbf{puede plantearse en cualquier pleito}, en la medida en que sea conducente para la resolución del conflicto. Lo ejerce cada juez en cualquier proceso en el que fue válidamente invocado, sin que exista un tribunal con una calificación exclusiva y excluyente para resolverlo, como sucede en otros órdenes jurídicos.
\end{definition}

En el sistema argentino, cualquier juez llamado a resolver un conflicto puede, en el marco de esa tarea y siempre que se trate de una cuestión justiciable que pueda resolver, ejercer el control de constitucionalidad y, como \textit{última ratio} (es decir, en el caso más extremo), \textbf{invalidar una norma} emanada del Congreso o del Poder Ejecutivo: por ejemplo, un decreto, un decreto de necesidad y urgencia o una ley.

\subsection{Comparación con los sistemas concentrados}

Declarar la invalidez de un acto o ley supone una eventual colisión entre poderes del Estado, porque es el Poder Judicial el que de algún modo desconoce una función legislativa o administrativa. Por ello, en otros regímenes esta facultad se supedita a determinados tribunales calificados.

En el sistema argentino, en cambio, cualquier juez, en el marco de un proceso que llegó a su conocimiento y en la medida en que sea conducente para resolver el conflicto bajo su jurisdicción, puede determinar la invalidez de una norma o disposición legislativa o administrativa. Esto no es posible, por ejemplo, en el sistema español. En el sistema continental europeo (por ejemplo, el italiano y el alemán) existen \textbf{tribunales constitucionales}: se ha optado por un tribunal, ya sea en sala o en salas, que lleva adelante el control de constitucionalidad.

\begin{table}[H]
\centering
\small
\begin{tabularx}{\textwidth}{
    >{\raggedright\arraybackslash}p{0.22\textwidth}
    >{\raggedright\arraybackslash}X
    >{\raggedright\arraybackslash}X
}
\toprule
\textbf{Característica} & \textbf{Sistema difuso (Argentina / EE.UU.)} & \textbf{Sistema concentrado (España / Alemania / Italia)} \\
\midrule
\textbf{¿Quién ejerce el control?} & Cualquier juez en cualquier proceso & Tribunal Constitucional especializado y exclusivo \\
\textbf{Acceso al justiciable} & Directo, en el mismo proceso judicial & Indirecto: se remite al tribunal especializado \\
\textbf{Alcance de la decisión} & Para el caso concreto (\textit{inter partes}) & Efecto general (\textit{erga omnes}) \\
\textbf{Cercanía geográfica} & Alta: el juez está en el lugar del conflicto & Menor: tribunal centralizado en la capital \\
\textbf{Instancias recursivas} & Múltiples instancias enriquecen el debate & Una única instancia constitucional \\
\textbf{Riesgo} & Sentencias contradictorias entre juzgados & Menor acceso para justiciables sin recursos \\
\textbf{Carácter} & Más democrático y accesible & Más especializado pero menos accesible \\
\bottomrule
\end{tabularx}
\end{table}

\section{Ventajas del sistema difuso}

Se sostiene una valoración favorable del sistema vigente: si bien tiene algunas imperfecciones y tiende a generar algunos marcos temporales de incertidumbre, resulta un sistema mucho más \textbf{democrático}, porque está al alcance de todos.

\begin{itemize}
    \item \textbf{Acceso de los abogados:} cualquier abogado, en la medida en que su capacidad de argumentación lo permita y tenga el caso para exponer, puede plantear firmemente la posibilidad de invalidez de una norma.
    \item \textbf{Cercanía y amplitud de los planteos:} el justiciable no debe esperar varios años a que un tribunal geográficamente lejano resuelva la cuestión y la reenvíe para ver si el juicio continúa; la cuestión se resuelve en el mismo proceso.
    \item \textbf{Instancias recursivas:} evitan que toda decisión se tome en un mismo tribunal sin instancia de revisión posterior. La revisión enriquece el discurso, porque puede llegar alguien que plantee algo que otros no plantearon, y además permite un tiempo de maduración de los temas: algunos se desarrollan con el tiempo y terminan de analizarse en su real dimensión pasado un tiempo.
\end{itemize}

\section{Costos del sistema difuso: la necesidad de un principio ordenador}

Hay que asumir un plazo de incertidumbre en el que puede haber:

\begin{itemize}
    \item sentencias contrapuestas;
    \item leyes generales puestas en duda o inaplicadas;
    \item enfrentamiento entre los poderes del Estado;
    \item la situación, a veces absurda pero real, de que una norma se aplique para algunos y para otros no.
\end{itemize}

Todas estas complejidades son los \textbf{costos de transacción} del sistema difuso.

Frente a todo lo valioso que tiene este sistema para lograr debates más complejos y meditados, no puede resignarse a que el sistema resulte caótico, ineficaz o contradictorio. Por ello, las decisiones constitucionales de la Corte deben tener un peso muy contundente y ser respetadas por los tribunales inferiores.

\begin{important}
El sistema difuso requiere de un \textbf{principio ordenador}. Sin él, caería en decisiones caóticas y contradictorias. Ese rol de ordenador lo cumple la Corte Suprema como intérprete máximo y definitivo de la Constitución.
\end{important}

\section{Cuadro Cornell: control de constitucionalidad}

\begin{longtable}{@{}>{\RaggedRight\arraybackslash}p{0.27\textwidth}>{\RaggedRight\arraybackslash}p{0.68\textwidth}@{}}
\toprule
\textbf{Preguntas / palabras clave} & \textbf{Notas / respuestas} \\
\midrule
\endfirsthead
\toprule
\textbf{Preguntas / palabras clave} & \textbf{Notas / respuestas} \\
\midrule
\endhead
\bottomrule
\endlastfoot

\textbf{¿Cuándo se ejerce un control de constitucionalidad?} &
Cuando, al interpretar la Constitución, se determina si existe una colisión entre normas inferiores y la Constitución, o cuál es el alcance de la propia norma fundacional. \\
\addlinespace[0.3em]\cmidrule[0.2pt]{1-2}\addlinespace[0.3em]

\textbf{¿Por qué se llama «difuso» al control de constitucionalidad argentino?} &
Porque puede plantearse en cualquier pleito, en la medida en que sea conducente para resolver el conflicto, y lo ejerce cada juez en cualquier proceso en el que fue válidamente invocado. No hay un tribunal con una calificación exclusiva y excluyente. Sigue muy fielmente la vertiente norteamericana. \\
\addlinespace[0.3em]\cmidrule[0.2pt]{1-2}\addlinespace[0.3em]

\textbf{¿Qué puede hacer un juez al ejercer el control de constitucionalidad?} &
Como última ratio, invalidar una norma emanada del Congreso o del Poder Ejecutivo (una ley, un decreto o un decreto de necesidad y urgencia), siempre que se trate de una cuestión justiciable que pueda resolver. \\
\addlinespace[0.3em]\cmidrule[0.2pt]{1-2}\addlinespace[0.3em]

\textbf{¿Qué diferencia hay entre el sistema difuso y el concentrado?} &
En el difuso cualquier juez ejerce el control en cualquier proceso; en el concentrado (por ejemplo, España, Italia y Alemania) el control lo lleva adelante un tribunal constitucional, en sala o en salas. Esto no es posible, por ejemplo, en el sistema español. \\
\addlinespace[0.3em]\cmidrule[0.2pt]{1-2}\addlinespace[0.3em]

\textbf{¿Por qué el sistema difuso es considerado más democrático?} &
Porque está al alcance de todos: cualquier abogado puede plantear la invalidez de una norma, el justiciable obtiene la resolución en el mismo proceso sin esperar a un tribunal lejano, y las instancias recursivas enriquecen el debate y permiten la maduración de los temas. \\
\addlinespace[0.3em]\cmidrule[0.2pt]{1-2}\addlinespace[0.3em]

\textbf{¿Cuáles son los costos del sistema difuso?} &
Un plazo de incertidumbre con sentencias contrapuestas, leyes generales puestas en duda o inaplicadas, enfrentamiento entre poderes y normas que se aplican para algunos y para otros no. Son los costos de transacción del sistema. \\
\addlinespace[0.3em]\cmidrule[0.2pt]{1-2}\addlinespace[0.3em]

\textbf{¿Qué necesita el sistema difuso para no resultar caótico?} &
Un principio ordenador: las decisiones constitucionales de la Corte deben tener un peso muy contundente y ser respetadas por los tribunales inferiores. \\
\addlinespace[0.3em]\cmidrule[0.2pt]{1-2}\addlinespace[0.3em]

\midrule
\multicolumn{2}{@{}p{0.95\textwidth}@{}}{%
\textbf{Resumen:} El control de constitucionalidad argentino es difuso: cualquier juez, en cualquier proceso, puede ejercerlo y, como última ratio, invalidar una norma. A diferencia de los sistemas concentrados con tribunal constitucional, resulta más democrático y accesible, pero tiene costos de transacción (incertidumbre y sentencias contradictorias). Por eso requiere un principio ordenador, función que cumple la Corte Suprema.} \\
\end{longtable}

%==============================================================
% CAPÍTULO 3
%==============================================================
\chapter{Los precedentes constitucionales}
\markboth{Capítulo \thechapter. Precedentes}{}

\section{La Corte Suprema como tribunal ordenador}

La decisión ordenadora, la que brinda la interpretación definitiva de la Constitución, es la de la \textbf{Corte}. Es la que define los precedentes constitucionales y la que funciona como organismo ordenador frente a las distintas sentencias que pueden ser contradictorias en los juzgados de primera instancia o en las cámaras.

En última instancia, los tribunales más elevados y superiores son los que armonizan las distintas decisiones: son la última cabeza del Poder Judicial, que jerárquicamente, por vía de recurso, termina dictando las directrices constitucionales.

\section{Obligatoriedad de los precedentes de la Corte}

El sistema argentino no es estrictamente un sistema de precedentes, aunque se asemeja cada vez más a uno. Aun manteniendo la idea de que el control de constitucionalidad tiene efectos para el caso, cuando se trata de decisiones de la Corte (o, dentro de una jurisdicción, de decisiones de un mismo tribunal superior, ya sea federal o local), se está ante verdaderas decisiones de carácter constitucional que \textbf{deben ser seguidas por los jueces inferiores}.

El sistema judicial argentino favorece que no haya un único tribunal constitucional, favorece la inmediatez entre quien acude a los tribunales y estos, y estimula las instancias recursivas en las que la cuestión constitucional puede ir variando y ser analizada. Todo este componente dialógico y democrático, que da acceso más directo a la posibilidad de ejercer un control de constitucionalidad en juicio, requiere sin embargo de un principio ordenador para no quedar en un \textit{loop} constante de desconocimiento de los principios de la Corte.

Las decisiones que tienen una densidad que permite proyectar el precedente a muchos casos similares deben ser seguidas obligatoriamente por los tribunales inferiores. Se abandona así la visión más reduccionista de la Corte de otras décadas, que mantenía a rajatabla el concepto de que los precedentes de la Corte solo eran obligatorios para el caso concreto.

\begin{important}
\textbf{Excepción al seguimiento del precedente.} Los tribunales inferiores solo podrían apartarse mediante una decisión meditada, muy prudente y muy justificada, si existiera alguna cuestión que no se planteó en el precedente constitucional, siempre justificando que hay algún aspecto que quizás no fue tenido en cuenta por la Corte Suprema.
\end{important}

\section{El overruling: cuando la Corte cambia su propio precedente}

\subsection{Concepto general}

La propia Corte también va modificando sus posiciones a lo largo del tiempo (\textit{overruling}: abandono del precedente). Incluso en los sistemas más ceñidos a la idea del \textit{common law} y del sistema de precedentes, siempre hay decisiones que pueden revertirse.

\subsection{Ejemplo histórico: la Corte Warren y el fin de la doctrina «separados pero iguales»}

En los años 60, la Corte americana, en la época de la lucha por los derechos civiles, modificó y revisó muchos de los precedentes que tenía vigentes desde décadas atrás, mucho más permeables y condescendientes con el sistema de segregacionismo. Fue la denominada Corte Warren la que desalojó esos precedentes.

La doctrina de \textbf{«separados pero iguales»} (\textit{separate but equal}) se vinculaba con los bienes, servicios y condiciones a los que podían acceder la población afroamericana y el resto de la población de los Estados Unidos. A partir de un precedente cercano al inicio del siglo XX, la Corte entendió que podía haber situaciones de segregacionismo siempre que no hubiera un servicio o una prestación diferente en calidad. Es decir, en la medida en que pudiera demostrarse que el acceso a determinados bienes y servicios era igual para toda la población, la discriminación o separación era admisible, porque no estaba en juego el principio de igualdad.

A fines de los años 50 y sobre todo en los 60, la Corte modificó completamente ese criterio, a partir de una nueva evaluación constitucional consustanciada con valores clave como la igualdad y la integración.
% TODO: contenido incompleto o ambiguo en las notas originales (la frase que describe el caso del servicio de salud separado y con los mismos recursos figura como «eso era válido», pero el desarrollo posterior indica que la Corte no lo consideró suficiente).
Aun cuando pudiera demostrarse que la población afroamericana tenía, por ejemplo, el mismo servicio de salud, por separado y con los mismos recursos, la Corte sostuvo que el solo hecho de que haya una segregación genera una dimensión clara de desigualdad en cuanto a la autopercepción de las minorías, que se ven desvalorizadas y separadas de actividades esenciales del resto de la población, y ello engendra y retroalimenta una idea de \textit{apartheid} y de segregacionismo. Por lo tanto, el concepto de «separados pero iguales» fue abandonado por la Corte.

Esto compromete una dimensión del análisis constitucional mucho más profunda, que recién se dio cuando la sociedad apoyó en algún punto esa hermenéutica jurídica y entendió que el segregacionismo por sí mismo, aun cuando pudiera haber iguales prestaciones, era negativo.

\subsection{Ejemplo doméstico: fallos Fayt y Schiffrin}

Un ejemplo argentino reciente es la modificación de la \textbf{doctrina Fayt}.

\paragraph{La doctrina Fayt.} Según esta doctrina, los jueces que cumplían 75 años podían permanecer en el cargo. Lo había establecido la Corte de los años 90, al entender que había existido una extralimitación de la Convención Constituyente al modificar el punto referido a la inamovilidad de los jueces. Esta doctrina se mantuvo durante mucho tiempo y llegó a generar la inaplicación de un punto de la Constitución: una suerte de control de constitucionalidad de la propia Constitución.

Se apoyaba en la vieja doctrina de la Corte según la cual las reformas constitucionales no pueden excederse de lo prefijado por la ley de reforma. En el fallo Fayt, que resulta muy discutible, además, la mayoría de la Corte que votó ese precedente estaba integrada por jueces que podían estar interesados en el resultado, por su propio interés en una futura inamovilidad en el cargo. Se resolvió que la modificación de la estabilidad superados los 75 años no estaba en la ley de reforma de la Constitución y que, por lo tanto, la reforma era inválida en ese punto. El precedente se mantuvo en el tiempo y permitió que muchos jueces que alcanzaban los 75 años se mantuvieran en el cargo.

\paragraph{El fallo Schiffrin.} Ante un caso que suscitó un recurso extraordinario, la misma Corte, ya con una nueva composición, volvió a analizar la cuestión y se apartó tajantemente de la doctrina de Fayt en la causa Schiffrin. La mayoría sostuvo que no hubo apartamiento de la ley que fijó la necesidad de reforma, porque entre los artículos incluidos estaban los referidos a la composición del Poder Judicial.

Aunque la ley preveía un núcleo de coincidencias básicas (una suerte de documento político que exteriorizaba el acuerdo al que habían llegado los partidos mayoritarios sobre la necesidad de reforma), ya preveía que podía reformarse el punto referido a la composición de los jueces del Poder Judicial, entre ellos la Corte. Con esa nueva mayoría, la Corte abandonó definitivamente el precedente Fayt y sostuvo con contundencia que la reforma constitucional era válida porque estaba contenida en la Constitución y que de ningún modo podía sostenerse que se hubiera excedido de la ley base, ya que su artículo referido al Poder Judicial estaba dentro de los que podían ser materia de modificación.

Dos de los jueces de la mayoría de Schiffrin habían sido constituyentes en el año 1994: Lorenzetti y Rosati.
% TODO: verificar grafía del apellido «Rosati» tal como figura en las notas originales.
Esto tiene un peso institucional muy fuerte.

\paragraph{El tiempo transcurrido desde Fayt.} Se destacó además, en relación con el voto de Rosati, el tiempo transcurrido desde Fayt: las presunciones de que la doctrina de Schiffrin podía implicar un ataque al Poder Judicial o una desestabilización de su independencia no se habían materializado en la práctica, de modo que esa preocupación quedó vacía de contenido. No había una manipulación de los jueces: el temor a una suerte de manoseo de la estabilidad y la independencia (siendo la estabilidad en el cargo una de las condiciones de la independencia) no estaba en juego.

\begin{table}[H]
\centering
\small
\begin{tabularx}{\textwidth}{
    >{\raggedright\arraybackslash}p{0.22\textwidth}
    >{\raggedright\arraybackslash}X
    >{\raggedright\arraybackslash}X
}
\toprule
\textbf{Aspecto} & \textbf{Fallo Fayt (1999)} & \textbf{Fallo Schiffrin (\textit{overruling})} \\
\midrule
\textbf{Cuestión debatida} & Validez de la reforma constitucional que fija retiro a los 75 años & Reexamen de la misma cuestión con nueva composición \\
\textbf{Posición adoptada} & La reforma era inválida: la Convención se extralimitó & La reforma era válida: estaba dentro de los artículos habilitados \\
\textbf{Mayoría} & Jueces con posible interés personal en el resultado & Dos integrantes habían sido constituyentes en 1994 \\
\textbf{Efectos prácticos} & Jueces mayores de 75 años podían continuar indefinidamente & Se restituyó la validez de la cláusula constitucional \\
\textbf{Impacto institucional} & Inaplicación de parte de la Constitución reformada & Corrección de un precedente fundado en premisas no verificadas \\
\bottomrule
\end{tabularx}
\end{table}

\section{Fallo Levinas: un precedente constitucional operativo}

En el precedente Levinas, muy reciente, la Corte, con una mayoría particular (eran tres los jueces que conformaban la mayoría), sentó el principio de que, en materia de derecho común, las sentencias de las \textbf{cámaras nacionales} son revisables por el Tribunal Superior de Justicia (TSJ) de la Ciudad como Tribunal Superior de la Causa, de forma previa a que la apelación vaya a revisión de la Corte. Las cámaras nacionales son aquellas que aplican e interpretan el derecho común en el ámbito de la Capital Federal: la Cámara Nacional en lo Civil, en lo Comercial, la Cámara Nacional en materia penal de delitos nacionales y en Derecho del Trabajo.

La Corte determinó así que, en materia ordinaria, más allá de que estos tribunales conforman la justicia nacional, los fallos de las cámaras nacionales pueden ser revisados por el TSJ de la Ciudad. Esta cuestión es una complejidad que todavía se está dilucidando en su real alcance: existe una primera sanción de la ley 27.802, en la que el Congreso sienta la idea de que debe haber un tránsito de la justicia nacional, en materia laboral, a la justicia eminentemente local.

\begin{note}
\textbf{Ley 27.802 --- Transferencia de la justicia nacional ordinaria a la Ciudad de Buenos Aires.} El Congreso de la Nación sancionó la ley 27.802, que establece el proceso de transferencia de la justicia nacional ordinaria con asiento en la Ciudad Autónoma de Buenos Aires a la jurisdicción local, en cumplimiento de los artículos 129 y de la cláusula transitoria decimoquinta de la Constitución Nacional reformada en 1994.
\end{note}

El fallo Levinas establece que, respecto de todas las sentencias que se notifiquen desde la notificación de ese fallo, las partes deberán interponer, si quieren impugnar, el recurso ante el TSJ. Es un ejemplo en el que la Corte, con una fisonomía mucho más operativa, da de forma muy contundente un \textbf{alcance general} a su precedente al fijar cuál es la jurisdicción revisora del derecho común en el ámbito de la Ciudad Autónoma de Buenos Aires.

\section{El sistema argentino de precedentes y el common law anglosajón}

Lo que diferencia en la práctica la visión más tradicional de los precedentes en el derecho americano es que hay un esfuerzo mucho más claro por dejar patente cuáles son los precedentes de la Corte que deben seguir los tribunales inferiores. En la práctica del precedente típico del derecho americano existe una ``gimnasia'' mucho más clara para que los tribunales inferiores identifiquen dónde hay una doctrina que deben seguir de forma obligatoria.

En cambio, en la costumbre jurídica argentina más reciente, la Corte ha dictado numerosas sentencias en las que destaca que los tribunales inferiores deben seguir estos precedentes constitucionales y que los jueces inferiores no pueden modificarlos y deben ceñirse a las hermenéuticas constitucionales establecidas por la Corte en fallos de proyección general. Lo ha hecho, sin embargo, como una labor más dispersa y quizás no tan contundente.

Por ejemplo, en el precedente Levinas la Corte hace un esfuerzo por darle una obligatoriedad hacia futuro muy clara, patente y operativa. De forma más rudimentaria, pero consistente, hoy la Corte argentina, más allá de las sucesivas integraciones que ha tenido, está haciendo un esfuerzo por darles una \textbf{densidad constitucional y obligatoria} a sus precedentes cuando así se requiera, para que se asiente la idea de que la Corte es el tribunal ordenador y de que, frente a decisiones constitucionales de proyección general, los tribunales inferiores deben someterse y seguir las directrices de sus fallos.

\begin{important}
Hay una tendencia a ir virando paulatinamente hacia un sistema que, si bien no es el típico de precedentes anglosajón, se le asemeja mucho: para evitar contradicciones y para evitar que las partes queden presas de la necesidad de apelar ante instancias innecesarias en sucesivas oportunidades, cuando hay un criterio constitucional claro y apoyado por la Corte, este tribunal está cumpliendo su misión ordenadora.
\end{important}

\section{Cuadro Cornell: precedentes constitucionales}

\begin{longtable}{@{}>{\RaggedRight\arraybackslash}p{0.27\textwidth}>{\RaggedRight\arraybackslash}p{0.68\textwidth}@{}}
\toprule
\textbf{Preguntas / palabras clave} & \textbf{Notas / respuestas} \\
\midrule
\endfirsthead
\toprule
\textbf{Preguntas / palabras clave} & \textbf{Notas / respuestas} \\
\midrule
\endhead
\bottomrule
\endlastfoot

\textbf{¿Qué rol cumple la Corte Suprema en el sistema difuso?} &
Es el tribunal ordenador y el intérprete máximo y definitivo de la Constitución. Define los precedentes constitucionales y armoniza las sentencias que pueden ser contradictorias en los juzgados de primera instancia o en las cámaras, como última cabeza del Poder Judicial. \\
\addlinespace[0.3em]\cmidrule[0.2pt]{1-2}\addlinespace[0.3em]

\textbf{¿Están obligados los jueces inferiores a seguir los precedentes de la Corte?} &
Sí, cuando se trata de decisiones con una densidad que permite proyectar el precedente a muchos casos similares. Se abandonó la visión reduccionista según la cual los precedentes solo obligaban para el caso concreto. Solo pueden apartarse mediante una decisión muy prudente y justificada, si hubiera una cuestión no considerada en el precedente. \\
\addlinespace[0.3em]\cmidrule[0.2pt]{1-2}\addlinespace[0.3em]

\textbf{¿Puede la Corte cambiar su propia doctrina?} &
Sí. Es el overruling o abandono del precedente, que se da incluso en los sistemas más ceñidos al common law. Ejemplos: la Corte Warren, que abandonó la doctrina «separados pero iguales», y el fallo Schiffrin, que abandonó la doctrina Fayt. \\
\addlinespace[0.3em]\cmidrule[0.2pt]{1-2}\addlinespace[0.3em]

\textbf{¿Por qué la Corte Warren abandonó la doctrina «separados pero iguales»?} &
Porque, con una nueva evaluación constitucional basada en valores como la igualdad y la integración, consideró que el solo hecho de la segregación genera una desigualdad clara, aun con prestaciones iguales, al desvalorizar a las minorías y retroalimentar una idea de apartheid. \\
\addlinespace[0.3em]\cmidrule[0.2pt]{1-2}\addlinespace[0.3em]

\textbf{¿Qué estableció el fallo Fayt y por qué es discutible?} &
Que la reforma constitucional que fija la estabilidad de los jueces hasta los 75 años era inválida porque la Convención Constituyente se había extralimitado respecto de la ley de reforma. Es discutible porque la mayoría estaba integrada por jueces que podían tener interés en el resultado y porque generó la inaplicación de un punto de la Constitución. \\
\addlinespace[0.3em]\cmidrule[0.2pt]{1-2}\addlinespace[0.3em]

\textbf{¿Qué estableció el fallo Schiffrin?} &
Que la reforma era válida, porque los artículos referidos al Poder Judicial estaban incluidos entre los que podían ser materia de modificación según la ley de reforma. Con una nueva composición de la Corte (dos de sus jueces habían sido constituyentes en 1994), se abandonó definitivamente el precedente Fayt. El temor a una afectación de la independencia judicial no se había materializado con el paso del tiempo. \\
\addlinespace[0.3em]\cmidrule[0.2pt]{1-2}\addlinespace[0.3em]

\textbf{¿Qué estableció el fallo Levinas y cómo obligó a los tribunales inferiores?} &
Que, en materia de derecho común, las sentencias de las cámaras nacionales son revisables por el TSJ de la Ciudad como Tribunal Superior de la Causa, antes de la apelación a la Corte. Para las sentencias notificadas desde la notificación del fallo, las partes deben interponer el recurso ante el TSJ. Es un ejemplo de precedente con alcance general y proyección hacia el futuro. \\
\addlinespace[0.3em]\cmidrule[0.2pt]{1-2}\addlinespace[0.3em]

\textbf{¿Se parece nuestro sistema al de precedentes anglosajón?} &
No es el típico sistema de precedentes, pero se le asemeja cada vez más. En el sistema americano hay una práctica clara para identificar los precedentes obligatorios; en la Argentina la Corte lo ha hecho de forma más dispersa, aunque hoy procura dar densidad constitucional y obligatoria a sus precedentes para cumplir su misión ordenadora. \\
\addlinespace[0.3em]\cmidrule[0.2pt]{1-2}\addlinespace[0.3em]

\midrule
\multicolumn{2}{@{}p{0.95\textwidth}@{}}{%
\textbf{Resumen:} La Corte Suprema es el tribunal ordenador del sistema difuso y sus decisiones constitucionales de proyección general obligan a los tribunales inferiores, que solo pueden apartarse con una justificación muy prudente. La Corte puede modificar sus propios precedentes (overruling), como ocurrió con la doctrina «separados pero iguales» y con el paso de Fayt a Schiffrin. El fallo Levinas ilustra un precedente operativo con alcance general. Todo ello indica un tránsito hacia una práctica cercana al sistema de precedentes anglosajón, con características propias.} \\
\end{longtable}

%==============================================================
% SÍNTESIS FINAL
%==============================================================
\chapter{Síntesis final}

La clase giró en torno a tres ejes articulados: la interpretación constitucional, el control de constitucionalidad y los precedentes judiciales.

En materia de interpretación, se analizaron el \textbf{textualismo}, que ciñe al juez al texto de la norma con el significado corriente de las palabras, y el \textbf{originalismo}, que agrega el anclaje histórico al momento de la sanción de la norma, destacando que la Corte argentina aplica una visión flexible del textualismo. El artículo 19 de la Constitución Nacional sirvió como ejemplo concreto de los dilemas hermenéuticos.

Luego se abordó el \textbf{sistema difuso} de control de constitucionalidad argentino, caracterizado por la posibilidad de que cualquier juez, en cualquier proceso, declare la invalidez de una norma. Este sistema fue comparado con el modelo concentrado europeo, destacándose su mayor democraticidad y accesibilidad, aunque también sus costos en términos de incertidumbre y potenciales contradicciones entre sentencias.

Finalmente, se analizó el papel de la \textbf{Corte Suprema} como intérprete máximo y tribunal ordenador, cuya función esencial es sentar precedentes constitucionales de seguimiento obligatorio para los tribunales inferiores. Se estudiaron ejemplos de overruling (el fallo Schiffrin, que revirtió el fallo Fayt sobre la edad de los jueces) y de precedentes operativos (el fallo Levinas, sobre la jurisdicción revisora en materia de derecho común en la Ciudad de Buenos Aires).

Todo el recorrido apunta a comprender que el sistema constitucional argentino está transitando hacia una práctica cada vez más cercana al sistema de precedentes anglosajón, aunque con características propias de nuestra tradición jurídica.

\end{document}
